\documentclass[11pt,a4paper]{article}

% ---------------------------------------------------------------
% Tipografia: Palatino, versalete, uma cor de destaque (bordô)
% ---------------------------------------------------------------
\usepackage[utf8]{inputenc}
\usepackage[T1]{fontenc}
\usepackage[brazil]{babel}
\usepackage[sc,osf]{mathpazo}
\usepackage{microtype}
\usepackage[top=2.4cm, bottom=2.6cm, left=2.6cm, right=2.6cm]{geometry}
\usepackage{amsmath}
\usepackage{amssymb}
\usepackage{xcolor}
\usepackage{tcolorbox}
\tcbuselibrary{skins, breakable}
\usepackage{enumitem}
\usepackage{booktabs}
\usepackage{array}
\usepackage{tabularx}
\usepackage{needspace}
\usepackage{qrcode}
\usepackage{hyperref}

\definecolor{vinho}{RGB}{122, 27, 42}
\definecolor{grafite}{RGB}{70, 70, 70}
\hypersetup{colorlinks=true, linkcolor=vinho, urlcolor=vinho}

\linespread{1.06}
\setlength{\parindent}{0pt}
\setlength{\parskip}{0.45em}

% ---------------------------------------------------------------
% Cabeçalhos
% ---------------------------------------------------------------
\newcommand{\secao}[1]{%
  \needspace{7\baselineskip}
  \vspace{1.8em}
  {\Large\bfseries #1}\\[-0.55em]
  {\color{black!30}\rule{\linewidth}{0.5pt}}
  \par\nobreak\smallskip
}

\newcommand{\objetivo}[1]{{\itshape\color{grafite}#1}\par\smallskip}

% ---------------------------------------------------------------
% Ficha de v\'ideo: n\'umero, t\'itulo, canal, dura\c{c}\~ao, link
% ---------------------------------------------------------------
\newcommand{\video}[5]{%
  \needspace{8\baselineskip}
  \par\addvspace{1.4em}
  {\large\bfseries\color{vinho}V\'ideo #1 \,\textemdash\, #2}\\[0.1em]
  {\small\scshape\color{grafite}#3 \,\textperiodcentered\, #4}\\[-0.5em]
  {\color{black!25}\rule{\linewidth}{0.4pt}}
  \par\nobreak\smallskip
  {\small\url{#5}}
  \par\smallskip
}

\newcommand{\campo}[2]{%
  \par{\small\scshape\color{vinho}#1}\ #2\par}

% ---------------------------------------------------------------
% Painéis: barra fina à esquerda
% ---------------------------------------------------------------
\newtcolorbox{painel}{blanker, breakable, left=12pt,
  borderline west={1.4pt}{0pt}{black!30},
  before skip=11pt, after skip=11pt}

\newtcolorbox{quadro}{blanker, breakable, left=12pt,
  borderline west={1.4pt}{0pt}{grafite},
  before skip=11pt, after skip=11pt,
  before upper={{\small\scshape\color{grafite}no quadro, depois do
    v\'ideo}\par\smallskip}}

\newtcolorbox{aviso}{blanker, breakable, left=12pt,
  borderline west={1.4pt}{0pt}{vinho},
  before skip=11pt, after skip=11pt,
  before upper={{\small\scshape\color{vinho}aten\c{c}\~ao}\par\smallskip}}

% ---------------------------------------------------------------
\begin{document}
% ---------------------------------------------------------------

% ---------------------------------------------------------------
% TÍTULO
% ---------------------------------------------------------------
\thispagestyle{empty}

{\scshape\color{vinho}cursinho solid\'ario \,\textperiodcentered\, f\'isica}\\[-0.5em]
{\color{vinho}\rule{\linewidth}{1.6pt}}\\[0.9em]
{\huge\bfseries Tr\^es V\'ideos}\\[0.35em]
{\LARGE Din\^amica --- as Leis de Newton em a\c{c}\~ao}\\[0.7em]
{\itshape\color{grafite}Situa\c{c}\~oes reais do dia a dia
\,\textperiodcentered\, 7 minutos de v\'ideo na aula inteira}
\hfill {\color{grafite}2026}\\[-0.3em]
{\color{black!30}\rule{\linewidth}{0.5pt}}

\vspace{1.0em}

\begin{painel}
\begin{itemize}[nosep, leftmargin=1.5em]
  \item S\~ao \textbf{filmagens de situa\c{c}\~oes reais}, n\~ao
        v\'ideo-aula: um teste de colis\~ao, um teste de freio e um
        experimento de propuls\~ao. A explica\c{c}\~ao continua sendo
        sua, no quadro.
  \item T\'itulo, canal e dura\c{c}\~ao foram conferidos.
  \item Os tr\^es juntos d\~ao \textbf{6 min 56 s}. O v\'ideo 1 sozinho
        cobre as tr\^es leis; se o tempo apertar de verdade, \'e o
        \'unico indispens\'avel.
  \item Assistir antes de levar, e conferir os links na v\'espera.
\end{itemize}
\end{painel}

\vspace{0.8em}

% ---------------------------------------------------------------
% MAPA
% ---------------------------------------------------------------
{\renewcommand{\arraystretch}{1.3}
\begin{tabularx}{\linewidth}{@{}lXcccl@{}}
  \multicolumn{6}{@{}l}{{\scshape\color{vinho}o que cada um cobre}}\\
  \addlinespace[2pt]
  \toprule
  & V\'ideo & 1\textordfeminine{} & 2\textordfeminine{} &
    3\textordfeminine{} & Bloco \\
  \midrule
  1 & Crash test: carro de 1959 contra carro de 2009
    & $\bullet$ & $\bullet$ & $\bullet$ & 0 ou 2 \\
  2 & Teste de freio: com ABS e sem ABS
    & $\bullet$ & $\bullet$ &           & 2 \\
  3 & Bal\~ao foguete preso a um fio
    &           & $\bullet$ & $\bullet$ & 3 \\
  \bottomrule
\end{tabularx}}

\vspace{0.6em}

\begin{aviso}
O v\'ideo 1 \'e o \'unico que mostra as tr\^es leis na mesma cena, e \'e
o mais curto dos tr\^es (1:13). Se a aula estiver atrasada, cortar o 2 e
o 3: eles refor\c{c}am e ligam com quest\~oes espec\'ificas da lista,
mas n\~ao introduzem nenhuma lei que o v\'ideo 1 j\'a n\~ao tenha
mostrado.
\end{aviso}

\newpage

% ===================================================================
\secao{V\'ideo 1 --- as tr\^es leis em uma colis\~ao}
% ===================================================================

\video{1}{1959 Chevrolet Bel Air vs.\ 2009 Chevrolet Malibu --- IIHS
crash test}{IIHS (Insurance Institute for Highway
Safety)}{1:13}{https://www.youtube.com/watch?v=C_r5UJrxcck}

\campo{quando}{Bloco 0 (abertura) \textbf{ou} Bloco 2, logo depois de
escrever $\vec{F}_R = m\vec{a}$. Se for usar v\'ideo em um momento s\'o
da aula, use aqui.}

\campo{o que \'e}{Dois carros de massa parecida batem de frente, um
contra o outro: um sedan de 1959, de a\c{c}o grosso e sem cinto de
tr\^es pontos, e um sedan de 2009. Sem narra\c{c}\~ao --- d\'a para
falar por cima e repetir o trecho em c\^amera lenta.}

\campo{pergunte antes}{``Qual carro voc\^es acham mais seguro: o antigo,
de a\c{c}o grosso, ou o moderno, que amassa todo?'' A turma quase sempre
escolhe o antigo. O v\'ideo mostra o contr\'ario: o Bel Air quase n\~ao
amassa e o boneco dentro dele \'e destru\'ido; o Malibu amassa muito e o
boneco sai bem.}

\campo{as tr\^es leis, na mesma cena}{}

\begin{enumerate}[leftmargin=2.4em, itemsep=0.55em]
  \item[\textbf{1\textordfeminine{}}] \textbf{O boneco continua.} O carro
        para contra o outro carro, mas o boneco segue para a frente com
        a velocidade que tinha --- ningu\'em o empurrou. Ele s\'o para
        quando o cinto, o airbag ou o volante o param. \'E o passageiro
        da freada do Bloco 0, em escala real.

  \item[\textbf{2\textordfeminine{}}] \textbf{Quem amassa protege.} Nos
        dois carros o boneco vai de $v$ a zero: a varia\c{c}\~ao de
        velocidade \'e a mesma, n\~ao d\'a para mudar. O que muda \'e o
        \textbf{tempo} em que isso acontece. A frente que amassa \'e uma
        frente que \textbf{estica o} $\Delta t$, e por isso derruba a
        for\c{c}a. A lataria r\'igida de 1959 n\~ao absorve nada: joga a
        desacelera\c{c}\~ao inteira, de uma vez, no corpo de quem est\'a
        dentro.

  \item[\textbf{3\textordfeminine{}}] \textbf{Nenhum dos dois ``ganha''
        a batida.} Os carros se empurram com for\c{c}as de
        \textbf{mesma intensidade e sentidos opostos}, sempre. O de 1959
        n\~ao empurra mais forte por ser mais r\'igido --- ele s\'o
        \textit{recebe} a mesma for\c{c}a de um jeito pior. \'E o
        argumento da Quest\~ao 9 (a m\~ae e o m\'ovel) aplicado a duas
        toneladas de a\c{c}o.
\end{enumerate}

\begin{quadro}
Escrever a 2\textordfeminine{} Lei na forma que explica o v\'ideo:
\[
  F = m\,a = m\,\frac{\Delta v}{\Delta t}
\]
e ao lado: \textbf{$\Delta v$ est\'a dado. S\'o d\'a para mexer no
$\Delta t$.} Deixar essa frase no quadro --- cinto, airbag, capacete,
luva de boxe e a lataria que amassa fazem todos a mesma coisa.

\smallskip
Para a 3\textordfeminine{} Lei, dois diagramas separados, um para cada
carro, com as duas setas de mesmo tamanho apontando para lados opostos.
\end{quadro}

\campo{prepara}{\textbf{Quest\~ao 5} (ENEM 2017, o gr\'afico dos cintos),
que ser\'a feita em sala: o cinto mais seguro \'e o de menor pico de
desacelera\c{c}\~ao, ou seja, o que espalha a frenagem no tempo. Depois
deste v\'ideo, a turma resolve a Q5 quase sem ajuda. Tamb\'em prepara a
\textbf{Quest\~ao 9}.}

\newpage

% ===================================================================
\secao{V\'ideo 2 --- por que o freio ABS para antes}
% ===================================================================

\video{2}{Teste dos freios ABS}{Quatro
Rodas}{2:02}{https://www.youtube.com/watch?v=x2BVJfaXH0A}

\campo{quando}{Bloco 2, item 4 (atrito), na hora de escrever
$\mu_e > \mu_c$.}

\campo{o que \'e}{Teste de revista com as duas frenagens lado a lado e a
dist\^ancia de parada medida: roda travada, derrapando, contra roda
girando, sem travar.}

\campo{pergunte antes}{``Para parar antes: travar a roda de vez, ou
n\~ao deixar travar?'' A intui\c{c}\~ao de quase todo mundo \'e pisar
com tudo e travar.}

\campo{as leis, no v\'ideo}{}

\begin{enumerate}[leftmargin=2.4em, itemsep=0.55em]
  \item[\textbf{1\textordfeminine{}}] \textbf{O carro n\~ao para
        sozinho.} Tirando o p\'e do acelerador ele seguiria em frente;
        quem muda o estado de movimento \'e uma for\c{c}a, e nesse caso
        a \'unica for\c{c}a que freia o carro \'e o \textbf{atrito dos
        pneus com o asfalto}. Nem \'e o freio: o freio segura a roda,
        quem segura o carro \'e o ch\~ao.

  \item[\textbf{2\textordfeminine{}}] \textbf{Mais atrito, mais
        desacelera\c{c}\~ao, menos dist\^ancia.} Como
        $a = F_{at}/m$, uma for\c{c}a de atrito maior d\'a uma
        desacelera\c{c}\~ao maior --- e o carro para em menos metros. As
        duas dist\^ancias medidas no v\'ideo s\~ao a 2\textordfeminine{}
        Lei com trena.
\end{enumerate}

\campo{o ponto que a turma erra}{Com a roda \textbf{travada} o pneu
\textbf{desliza} sobre o asfalto: vale o atrito \textbf{cin\'etico}. Com
a roda \textbf{girando} o ponto do pneu que toca o ch\~ao \textbf{n\~ao
desliza}: vale o atrito \textbf{est\'atico}, que \'e maior. Da\'i o ABS
parar antes. E o atrito \textbf{n\~ao depende da \'area de contato} ---
pneu mais largo n\~ao entra nessa conta.}

\begin{quadro}
\[
  f_{at} = \mu N \qquad\text{com}\qquad \mu_e > \mu_c
\]
Embaixo: \textbf{roda travada $\Rightarrow$ desliza $\Rightarrow$
cin\'etico. Roda girando $\Rightarrow$ n\~ao desliza $\Rightarrow$
est\'atico.}
\end{quadro}

\campo{prepara}{\textbf{Quest\~ao 11} (ENEM PPL 2012), que \'e este
v\'ideo em forma de enunciado --- inclusive as alternativas erradas
sobre \'area de contato.}

% ===================================================================
\secao{V\'ideo 3 --- a\c{c}\~ao e rea\c{c}\~ao sem ter em que se apoiar}
% ===================================================================

\video{3}{Bal\~ao Foguete e a 3\textordfeminine{} Lei de Newton}{PET
F\'isica UFAM}{3:41}{https://www.youtube.com/watch?v=R1W2Nohk8Yg}

\campo{quando}{Bloco 3, no exemplo do foguete.}

\campo{o que \'e}{Um bal\~ao cheio de ar, preso a um canudo que corre por
um fio esticado. Solta o bico e ele dispara. D\'a para \textbf{fazer ao
vivo} com bal\~ao, canudo, fita e linha --- nesse caso o v\'ideo serve
de roteiro de montagem, ou de plano B.}

\campo{pergunte antes}{``O bal\~ao anda porque o ar que sai
\textbf{empurra o ar da sala}?'' \'E o que quase todo mundo acha, e \'e
a mesma ideia errada da alternativa (e) da Quest\~ao 10, do astronauta
que ``nadaria'' no v\'acuo.}

\campo{as leis, no v\'ideo}{}

\begin{enumerate}[leftmargin=2.4em, itemsep=0.55em]
  \item[\textbf{3\textordfeminine{}}] \textbf{O par est\'a em corpos
        diferentes.} O bal\~ao empurra o ar para tr\'as; o ar empurra o
        bal\~ao para a frente. S\~ao duas for\c{c}as de mesma
        intensidade em \textbf{corpos diferentes} --- por isso n\~ao se
        cancelam e o bal\~ao acelera. E o que empurra o bal\~ao \'e o ar
        que \textbf{saiu dele}, n\~ao o ar de fora: por isso o mesmo
        princ\'ipio funciona no v\'acuo.

  \item[\textbf{2\textordfeminine{}}] \textbf{Enquanto sai ar, ele
        acelera.} A propuls\~ao dura s\'o enquanto h\'a ar saindo;
        acabou o ar, acabou a for\c{c}a, e o bal\~ao segue perdendo
        velocidade s\'o pelo atrito. \'E a diferen\c{c}a entre
        \textbf{ter for\c{c}a} e \textbf{ter velocidade} --- a mesma
        confus\~ao da Quest\~ao 3.
\end{enumerate}

\begin{quadro}
Dois diagramas lado a lado, um para o bal\~ao e um para o ar expelido,
com as setas de mesmo tamanho e sentidos opostos. Do lado, escrever:
\textbf{n\~ao precisa de apoio no ar de fora}.
\end{quadro}

\campo{prepara}{\textbf{Quest\~ao 10} (UTFPR, astronauta que se move
expelindo g\'as) e o argumento da \textbf{Quest\~ao 9}.}

\newpage

% ===================================================================
\secao{P\'agina para projetar no fim da aula}
% ===================================================================

\objetivo{Projetar no Bloco 6, junto com o aviso da tarefa, para quem
quiser rever em casa. Os alunos apontam a c\^amera do celular para o
c\'odigo.}

\vspace{1.5em}

\begin{center}
{\renewcommand{\arraystretch}{1.4}
\begin{tabular}{@{}p{3.7cm}p{9.0cm}@{}}
  \qrcode[height=2.8cm,padding]{https://www.youtube.com/watch?v=C_r5UJrxcck} &
  \raisebox{1.0cm}{\parbox[t]{9.0cm}{%
    \textbf{Crash test: carro de 1959 contra carro de 2009}\\
    {\small\scshape\color{grafite}iihs
    \,\textperiodcentered\, 1:13}\\
    {\small As tr\^es leis na mesma colis\~ao. Repare no boneco, n\~ao
    na lataria.}}}
  \\[1.8em]
  \qrcode[height=2.8cm,padding]{https://www.youtube.com/watch?v=x2BVJfaXH0A} &
  \raisebox{1.0cm}{\parbox[t]{9.0cm}{%
    \textbf{Teste dos freios ABS}\\
    {\small\scshape\color{grafite}quatro rodas
    \,\textperiodcentered\, 2:02}\\
    {\small Por que n\~ao travar a roda faz o carro parar antes.}}}
  \\[1.8em]
  \qrcode[height=2.8cm,padding]{https://www.youtube.com/watch?v=R1W2Nohk8Yg} &
  \raisebox{1.0cm}{\parbox[t]{9.0cm}{%
    \textbf{Bal\~ao foguete e a 3\textordfeminine{} Lei}\\
    {\small\scshape\color{grafite}pet f\'isica ufam
    \,\textperiodcentered\, 3:41}\\
    {\small D\'a para montar em casa com bal\~ao, canudo, fita e
    linha.}}}
\end{tabular}}
\end{center}

\end{document}
